% Глава 2. Обзор существующих методов

\section{Обзор существующих методов}
\label{sec:literature_review}

\subsection{Диффузионные модели генерации изображений}

Диффузионные модели (diffusion models) образуют класс генеративных моделей, основанных на постепенном зашумлении данных и последующем обучении нейронной сети обращать этот процесс. Класс стал одним из ключевых в синтезе изображений, превзойдя генеративно-состязательные сети (generative adversarial networks, GAN)~\cite{dhariwal2021diffbeatgans} по основным показателям качества и разнообразия генерируемых изображений.

Основополагающей для современной эры диффузионных моделей стала работа Хо и соавт.~\cite{ho2020ddpm}, предложивших модель денойзингового диффузионного вероятностного процесса (Denoising Diffusion Probabilistic Model, DDPM). Подход определяет прямой марковский процесс (forward Markov process), в ходе которого к исходному изображению на протяжении $T$ шагов последовательно добавляется гауссовский шум, постепенно превращая его в чистый нормальный шум. Формально, для шага $t$ прямой процесс задаётся как:

\begin{equation}
    q(x_t \mid x_{t-1}) = \mathcal{N}(x_t;\, \sqrt{1-\beta_t}\, x_{t-1},\, \beta_t \mathbf{I})
    \label{eq:ddpm_forward}
\end{equation}

где параметры $\beta_t$ задают расписание шумования (noise schedule). Обратный процесс, реализуемый нейронной сетью, обучается предсказывать добавленный на каждом шаге шум посредством минимизации взвешенного среднеквадратичного отклонения (так называемая $\varepsilon$-параметризация), что эквивалентно денойзинговому согласованию оценок (denoising score matching) при соответствующей перепараметризации~\cite{song2021scoresde}. В качестве архитектуры обратного денойзера использовалась свёрточная сеть в форме U-Net с механизмами внимания. На наборе данных CIFAR-10 модель DDPM достигла показателя FID (Fréchet Inception Distance) 3,17~--- на момент публикации один из наилучших результатов среди likelihood-based генеративных моделей. Работа~\cite{ho2020ddpm} заложила концептуальную и экспериментальную основу всех последующих разработок.

Главное практическое ограничение DDPM~--- высокая вычислительная стоимость генерации: на один образец требуются сотни шагов денойзинга. Сонг и соавт.~\cite{song2021ddim} устранили это ограничение, предложив детерминированный процесс выборки на основе немарковских цепей~--- модель DDIM (Denoising Diffusion Implicit Model). DDIM ускоряет генерацию в 10--50 раз по сравнению с исходным DDPM при сопоставимом качестве. Детерминированность процесса открывает интерполяцию в пространстве скрытых представлений, необходимую для задач редактирования и переноса стиля.

Ромбах и соавт.~\cite{rombach2022ldm} перенесли диффузионный процесс из пространства пикселей в пространство скрытых переменных (latent space) и предложили модель латентной диффузии (Latent Diffusion Model, LDM). Архитектурно LDM состоит из трёх компонентов: вариационного автоэнкодера (Variational Autoencoder, VAE), сжимающего изображение в компактное скрытое пространство; диффузионной модели, работающей непосредственно в этом пространстве; и механизма условного управления через перекрёстное внимание (cross-attention), реализующего текстовое обусловливание на основе языковых эмбеддингов. Это резко снизило вычислительную стоимость обучения и инференса при сохранении качества генерации. На базе архитектуры LDM разработана Stable Diffusion --- открытая модель, ставшая де-факто стандартной платформой для методов управления генерацией, персонализации и переноса стиля. Метод B-LoRA~\cite{frenkel2024blora} изначально разработан на SDXL~--- расширенной версии Stable Diffusion (адаптация к FLUX.1 рассмотрена в разделе~\ref{sec:method}).

Параллельно с развитием диффузионных подходов активно исследовались методы переноса стиля на основе свёрточных нейронных сетей. Работа Гатиса и соавт.~\cite{gatys2016nst} установила формальное разграничение между стилем и содержанием изображения: глубокие признаки VGG-сети, извлечённые из разных слоёв, описывают содержание изображения, тогда как статистика корреляций этих признаков, формализованная через матрицу Грама (Gram matrix), кодирует стилистические характеристики. Метод итеративной оптимизации позволял комбинировать стиль одного изображения с содержанием другого. Несмотря на вычислительную затратность и отсутствие управления через текстовый запрос, этот подход заложил концептуальную основу всех последующих исследований по разделению стиля и содержания, включая B-LoRA.

К 2023 г. диффузионные модели перешли от пиксельного пространства к латентному и от марковских цепей DDPM к детерминированным семплерам DDIM, оформившись как универсальная платформа для синтеза и редактирования изображений с текстовым и визуальным обусловливанием.


\subsection{Диффузионные трансформеры}

По мере роста масштаба диффузионных моделей архитектура U-Net, основанная на свёрточных блоках с остаточными связями и механизмами пространственного внимания, начала проявлять структурные ограничения. Во-первых, свёрточные операции обладают локальным рецептивным полем (receptive field), что затрудняет моделирование дальних пространственных зависимостей. Во-вторых, масштабирование U-Net по числу параметров происходит неоднородно и требует ручной настройки архитектурных гиперпараметров. В-третьих, режим условного обусловливания через перекрёстное внимание в U-Net реализован отдельно от механизма самовнимания, что потенциально ограничивает качество согласования текстового и визуального потоков. Эти ограничения создали предпосылки для перехода к архитектурам на основе трансформеров (transformer-based architectures).

Теоретическую основу всех современных архитектур трансформеров составляет механизм многоголового самовнимания (multi-head self-attention), предложенный в работе Васвани и соавт.~\cite{vaswani2017attention}. Механизм внимания вычисляет взвешенные суммы значений (values) на основе попарного сходства запросов (queries) и ключей (keys):

\begin{equation}
    \text{Attention}(Q, K, V) = \text{softmax}\!\left(\frac{QK^\top}{\sqrt{d_k}}\right) V
    \label{eq:attention}
\end{equation}

Трансформеры устойчиво масштабируются с ростом числа параметров и вычислительных ресурсов (законы масштабирования, scaling laws), что обусловило их выбор как альтернативы свёрточным архитектурам в роли денойзера в диффузионных моделях.

Первое систематическое исследование трансформеров в роли денойзеров диффузионных моделей выполнено Пиблзом и Се~\cite{peebles2023dit} с архитектурой DiT (Diffusion Transformer). Ключевая идея DiT состоит в разбиении скрытого представления на патчи (patches) --- небольшие прямоугольные участки --- и их обработке стандартными блоками трансформера с механизмом самовнимания. Сигнал условного управления, включающий временной шаг диффузии и метку класса, вводится посредством слоёв адаптивной нормализации (adaptive layer normalization, adaLN) с инициализацией нулём (adaLN-Zero), что обеспечивает нулевое воздействие условного сигнала на выход блока в начале обучения и существенно стабилизирует процесс тренировки. На наборе данных ImageNet с разрешением 256×256 пикселей модель DiT достигла значения FID равного 2,27, превзойдя существующие диффузионные модели на основе U-Net. Принципиальный результат: качество модели предсказуемо улучшается с ростом числа параметров и вычислительного бюджета (FLOP), что подтверждает выполнение законов масштабирования для этой архитектуры~\cite{peebles2023dit}.

Эссер и соавт.~\cite{esser2024sd3} разработали мультимодальный DiT (Multimodal DiT, MM-DiT), легший в основу модели Stable Diffusion~3. В отличие от классического DiT, обрабатывающего единый поток токенов, MM-DiT использует архитектуру с двумя параллельными потоками --- отдельными для токенов изображения и токенов текста. Эти потоки взаимодействуют посредством совместного механизма внимания (joint attention), в котором матрицы запросов, ключей и значений для обоих потоков объединяются перед вычислением внимания. Такая архитектура сохраняет за каждым потоком собственные специализированные веса и одновременно обеспечивает двунаправленное взаимодействие визуальных и текстовых представлений. Дополнительно в~\cite{esser2024sd3} предложено улучшенное расписание шумования на основе выпрямленного потока (rectified flow). Совокупность архитектурных решений MM-DiT задала новый уровень качества генерации изображений по текстовому описанию.

Крупнейшая на 2025~г. открытая трансформерная архитектура для генерации изображений~--- FLUX.1, разработанная компанией Black Forest Labs~\cite{blackforest2024flux}. FLUX.1~--- гибридная архитектура, совмещающая блоки MM-DiT с параллельным механизмом внимания (parallel attention) при 12~миллиардах параметров. В качестве текстовых энкодеров используются T5-XXL и CLIP~\cite{radford2021clip}. Архитектура FLUX.1 включает два типа блоков: блоки двойного потока (double stream blocks), реализующие принцип MM-DiT с отдельными потоками для изображений и текста, и блоки единого потока (single stream blocks), обрабатывающие объединённое представление на поздних этапах обработки.

FLUX.1 распространяется в трёх вариантах: \texttt{[pro]}~--- проприетарная версия с наилучшим качеством; \texttt{[dev]}~--- версия с открытыми весами для некоммерческого использования; \texttt{[schnell]}~--- дистиллированная версия под Apache~2.0, генерирующая изображения за 4 шага инференса. По пользовательским предпочтениям FLUX.1 превосходит конкурирующие системы на большинстве оценочных бенчмарков~\cite{blackforest2024flux}.

Переход от U-Net к трансформерным архитектурам меняет организацию информационных потоков: от локальных свёрточных операций~--- к глобальному вниманию, от раздельных потоков текста и изображения~--- к их совместной обработке. Архитектурное разделение на блоки двойного и единого потока в FLUX.1 определяет постановку метода блочно-специфической адаптации в разделе~\ref{sec:method}.


\subsection{Выпрямленные потоки}

Параллельно с развитием архитектур денойзеров улучшался и формализм построения генеративного процесса. Диффузионные модели на марковских цепях имеют ряд недостатков: траектории от шума к данным криволинейны, что требует большого числа шагов для их аппроксимации с приемлемой точностью. Подход на основе потоков (flow-based methods) предлагает альтернативу.

В работе Липмана и соавт.~\cite{lipman2023flow} был предложен метод согласования потока (Flow Matching, FM) --- общий фреймворк для обучения непрерывных нормализующих потоков (Continuous Normalizing Flows, CNF) без использования явного моделирования ОДУ в процессе обучения. Ключевая идея заключается в задании условного вектора скорости, переносящего образцы от шума к данным по оптимальным траекториям. При использовании интерполяции оптимальным транспортом (Optimal Transport, OT) данные траектории оказываются прямолинейными.

Концептуально близкий, но независимо разработанный подход был предложен в работе Лю и соавт.~\cite{liu2022rectifiedflow} под названием выпрямленный поток (Rectified Flow). Выпрямленный поток транспортирует вероятностные меры вдоль прямолинейных траекторий посредством обыкновенных дифференциальных уравнений. Формально, процесс интерполяции задаётся как:

\begin{equation}
    X_t = (1 - t)\, X_0 + t\, X_1, \quad t \in [0, 1]
    \label{eq:rectflow_interp}
\end{equation}

где в обозначениях оригинальной статьи~\cite{liu2022rectifiedflow} $X_0$~--- данные, $X_1$~--- нормальный шум, а нейронная сеть обучается предсказывать постоянную скорость перехода $v = X_1 - X_0$ (в разделе~\ref{sec:method} принята обратная конвенция \texttt{diffusers}~--- содержательный смысл не меняется). Ключевым преимуществом прямолинейных траекторий является возможность использования более редкого шага интегрирования при инференсе: чем прямолинейнее траектория, тем меньше шагов требуется для достаточно точного её следования.

Дополнительно в~\cite{liu2022rectifiedflow} предложена процедура рекурсивного выпрямления (rectification): путём повторного применения операции Reflow на основе ранее сгенерированных пар (данные, шум) можно дополнительно выпрямить траектории, в пределе приближаясь к одношаговой генерации. Формализм выпрямленного потока используется в архитектурах MM-DiT и FLUX.1 в качестве генеративного транспорта. Вариант FLUX.1~\texttt{[schnell]}, способный генерировать изображения за 4~шага инференса, получен из FLUX.1~\texttt{[dev]} не через процедуру рекурсивного Reflow, а через дистилляцию с состязательной целью (adversarial distillation)~\cite{blackforest2024flux}.


\subsection{Методы эффективного дообучения}

Полное дообучение современных диффузионных моделей в миллиарды параметров вычислительно дорого. Для задач персонализации, адаптации стиля и предметно-ориентированной генерации нужны методы, обновляющие лишь небольшую часть весов. Этот класс подходов известен как методы эффективного дообучения с малым числом параметров (Parameter-Efficient Fine-Tuning, PEFT).

Центральный метод этой области~--- низкоранговая адаптация (Low-Rank Adaptation, LoRA), предложенная Ху и соавт.~\cite{hu2022lora}. Ключевая идея: изменения весов предобученной модели при адаптации к новой задаче обладают низким эффективным рангом. Формально, вместо обновления матрицы весов $W \in \mathbb{R}^{d \times k}$ целиком, LoRA параметризует приращение весов как произведение двух матриц малого ранга:

\begin{equation}
    \Delta W = B A, \quad B \in \mathbb{R}^{d \times r},\; A \in \mathbb{R}^{r \times k}, \quad r \ll \min(d, k)
    \label{eq:lora_lit}
\end{equation}

При инференсе приращение $\Delta W$ масштабируется гиперпараметром $\alpha / r$ и складывается с замороженными предобученными весами; при выполнении слияния весов LoRA-адаптера и базовой модели (merge) дополнительная задержка инференса не возникает, тогда как при <<горячем>> подключении адаптера сохраняется небольшое снижение пропускной способности~\cite{hu2022lora}. На практике в диффузионных моделях LoRA обычно применяется к матрицам запросов ($Q$), ключей ($K$), значений ($V$) и выходной проекции в слоях внимания. При ранге $r \in \{4, 8\}$ число обучаемых параметров сокращается в 10~000 раз по сравнению с полным дообучением для моделей масштаба GPT-3 без ощутимой потери качества адаптации~\cite{hu2022lora}.

Из расширений базового подхода выделяется DoRA (Weight-Decomposed Low-Rank Adaptation)~\cite{liu2024dora}, декомпозирующий матрицу весов на компоненты величины (magnitude) и направления (direction) и применяющий LoRA только к компоненте направления, что точнее воспроизводит паттерны полного дообучения. Систематическое сравнение различных вариантов LoRA --- LoHa, LoKr, OFT, BOFT, DoRA --- предпринято авторами библиотеки LyCORIS~\cite{yeh2024lycoris}, предоставляющей комплексный фреймворк оценки применительно к диффузионным моделям.

Применение LoRA к диффузионным моделям для задач персонализации является прямым предшественником метода B-LoRA: если базовый LoRA обрабатывает все блоки модели единообразно, то B-LoRA совершает принципиальный шаг --- ограничивает обучение двумя специфическими блоками, опираясь на гипотезу о структурной специализации блоков применительно к стилю и содержанию.

\begin{figure}[h]
    \centering
    \includegraphics[width=0.98\textwidth]{diffusion_process.png}
    \caption{Концептуальная схема диффузионной генерации и точка
             подключения LoRA-адаптера. Сверху~--- прямой марковский
             процесс зашумления (используется только при обучении);
             снизу~--- обратный процесс, реализуемый нейросетью-денойзером
             ($\varepsilon$-параметризация в DDPM, поле скоростей $v_\theta$
             в формализме rectified flow). Денойзер принимает три условных
             входа: зашумлённый сигнал $x_t$, шаг диффузии $t$ и текстовый
             эмбеддинг. Архитектура денойзера менялась от U-Net (DDPM,
             Stable~Diffusion) к диффузионным трансформерам DiT и MM-DiT
             (FLUX.1, Stable~Diffusion~3). LoRA-адаптация добавляет
             низкоранговые обновления $\Delta W = BA$ к замороженным
             линейным проекциям внимания, не модифицируя сам денойзер.}
    \label{fig:diffusion_process}
\end{figure}


\subsection{Методы персонализации и переноса стиля}

Персонализация генеративных моделей подразумевает адаптацию предобученной модели к воспроизведению специфического визуального концепта по небольшому числу референсных изображений. Перенос стиля~--- частный случай, при котором стилистические характеристики референсного изображения применяются к новому контенту с сохранением его семантики.

Один из первых систематических подходов к персонализации~--- текстовое инвертирование (Textual Inversion)~\cite{gal2023textualinversion}. Метод обучает вектор встраивания для нового псевдослова по 3--5 референсным изображениям при замороженных весах модели. Ключевое достоинство~--- полное сохранение генеративных возможностей базовой модели; ограничение~--- относительно слабая способность воспроизводить детали.

Более мощный подход~--- DreamBooth~\cite{ruiz2023dreambooth}, выполняющий полное дообучение модели на 3--5 изображениях с редким идентификатором-токеном. Для предотвращения семантического дрейфа используется потеря сохранения априорного распределения (prior preservation loss). DreamBooth даёт высокое качество воспроизведения субъекта ценой дорогостоящего полного дообучения.

Иной подход реализован в IP-Adapter~\cite{ye2023ipadapter}: отдельный механизм перекрёстного внимания для признаков изображения-условия обучается при замороженной базовой модели. Адаптер с 22 миллионами параметров проецирует признаки референсного изображения в то же пространство ключей и значений, что и текст, сохраняя полную управляемость через текстовый запрос.

Метод StyleAligned~\cite{hertz2024stylealigned} обеспечивает стилистическую согласованность между несколькими генерируемыми изображениями без обучения: статистики ключей и значений механизма самовнимания разделяются между всеми изображениями батча; применяется адаптивная нормализация по экземпляру (Adaptive Instance Normalization, AdaIN)~\cite{huang2017adain}.

Метод InstantStyle~\cite{wang2024instantstyle} предлагает два решения для улучшения переноса стиля: вычитание признаков содержания из признаков стиля в пространстве IP-Adapter и инъекцию стиля исключительно в блоки перекрёстного внимания, специфичные для стиля. InstantStyle опирается на предположение о блочной специализации, разделяемое методом B-LoRA.

Метод DEADiff~\cite{qi2024deadiff} реализует разделённый подход к кодированию стиля и содержания на архитектуре Q-Former с отдельными токенами-инструкциями. Извлечённые признаки инъецируются в различные слои перекрёстного внимания денойзера. Принципиальная черта~--- использование пар изображений с одинаковым стилем, но разным содержанием, что предотвращает утечку информации между потоками.

Метод B-LoRA~\cite{frenkel2024blora}~--- непосредственный предшественник предлагаемого здесь подхода. Опубликованный на ECCV 2024, B-LoRA строится на эмпирическом наблюдении: в архитектуре SDXL разные блоки кодируют разные аспекты визуальной информации. Идентификация целевых блоков выполнена авторами процедурой инъекции текстовых эмбеддингов (\textit{prompt injection}): подстановка стилевого либо содержательного промпта в каждый блок и анализ отклика генерации позволили установить, что блок \texttt{down\_blocks[2].attentions[1]} в иерархии Down-энкодера (обозначаемый в работе как \textit{W1}) преимущественно отвечает за семантическое содержание изображения, тогда как блок \texttt{up\_blocks[0].attentions[1]} в иерархии Up-декодера (обозначаемый как \textit{W2}) кодирует стилистические характеристики. На основании этого предлагается обучать LoRA-адаптеры исключительно на двух названных блоках, используя в качестве обучающих данных одно-единственное изображение.

Принципиальный результат такого ограниченного обучения~--- неявное разделение стиля и содержания (implicit style/content disentanglement) без явного надзорного сигнала. Обученные адаптеры распадаются на два типа: B-LoRA содержания (Content B-LoRA), ассоциированный с блоком~4, и B-LoRA стиля (Style B-LoRA), ассоциированный с блоком~5. Данные адаптеры могут применяться независимо или совместно, обеспечивая несколько режимов работы: перенос стиля при новом контенте; текстово-управляемая стилизация; генерация изображений, согласованных по стилю. При этом никаких дополнительных шагов обучения для инференса не требуется.

Обучение B-LoRA производится на одном изображении и занимает порядка пяти минут на ускорителе уровня NVIDIA~A100 при числе итераций около~1\,000~\cite{frenkel2024blora}. Потеря обучения является стандартной диффузионной потерей денойзинга применительно к двум целевым блокам. Переобучения не наблюдается благодаря двум факторам: крайне малому числу обучаемых параметров и регуляризирующему эффекту замороженных весов базовой модели, обеспечивающему разнообразие при инференсе.

Экспериментальные результаты~\cite{frenkel2024blora} демонстрируют, что B-LoRA существенно превосходит Textual Inversion~\cite{gal2023textualinversion} и DreamBooth~\cite{ruiz2023dreambooth} по метрикам стилистической верности применительно к задаче переноса стиля в рамках SDXL. Однако B-LoRA имеет принципиальное ограничение: метод разработан и валидирован исключительно на архитектуре SDXL на основе U-Net. Применимость базовой идеи B-LoRA к диффузионным трансформерам~--- центральная проблема настоящей работы.

Магистральное направление развития персонализации и переноса стиля~--- методы, опирающиеся на структурную специализацию блоков модели; B-LoRA наиболее последовательно воплощает эту тенденцию.


\subsection{Адаптация методов переноса стиля к архитектуре FLUX}

С ростом популярности FLUX.1 как генеративной платформы ряд исследований поставил задачу адаптации существующих методов переноса стиля к её специфической архитектуре.

Хуанг и соавт.~\cite{huang2024iclora} демонстрируют способность FLUX.1 к генерации в контексте нескольких изображений (in-context generation): конкатенация нескольких изображений в единое контекстное окно активирует латентные возможности модели, которые затем дообучаются стилезависимым LoRA малого ранга (20--100 образцов). Работа устанавливает совместимость LoRA и FLUX.1, но не исследует специализацию отдельных блоков.

Метод FluxSpace~\cite{dalva2024fluxspace} исследует структуру пространства внимания FLUX.1 с точки зрения редактирования. Авторы показывают, что выходы блоков совместного внимания трансформера образуют семантически значимое пространство представлений, пригодное для управляемого редактирования без обучения на двух уровнях~--- грубой стилистики и тонких атрибутов. Это косвенно свидетельствует о специализации разных блоков FLUX.1 на разные аспекты визуальной информации~--- предположение, лежащее в основе адаптации идей B-LoRA к FLUX.

Наиболее близкая к задаче настоящей работы~--- SplitFlux~\cite{yang2025splitflux}, в которой авторы переносят идею блочной декомпозиции LoRA на архитектуру FLUX.1. Формально следуя гипотезе блочной специализации, SplitFlux расширяет её двумя дополнительными модулями. Авторы устанавливают, что среди блоков единого потока FLUX.1 ранние блоки преимущественно кодируют структурное содержание, тогда как поздние блоки несут ответственность за стилистические характеристики; на основе этого разделения предлагается раздельное обучение LoRA-адаптеров. Помимо базового разделения, SplitFlux вводит два дополнительных механизма: адаптацию с ограничением ранга (Rank-Constrained Adaptation, RCA) для предотвращения деградации идентичности субъекта и визуально-управляемый LoRA (Visual-Gated LoRA Adaptation, VGRA) для фокусировки адаптации на значимых областях.

У SplitFlux есть ряд ограничений. Введение RCA и VGRA усложняет систему по сравнению с концептуальной простотой исходного B-LoRA. Авторы не проводят систематического аблационного анализа роли отдельных блоков. Задействованы только блоки единого потока, тогда как роль блоков двойного потока остаётся неисследованной. Комплексная архитектура затрудняет интерпретацию вклада отдельных компонентов.

Анализ рассмотренных работ выявляет четыре лакуны: отсутствие чистой адаптации B-LoRA к FLUX.1 без дополнительных компонентов; отсутствие систематического анализа блочной специализации FLUX.1; неисследованность роли блоков двойного потока; отсутствие верификации применимости гипотезы B-LoRA при переходе от U-Net к DiT-архитектуре.


\subsection{Выводы по обзору}

Проведённый обзор литературы охватывает ключевые направления исследований по диффузионным моделям, трансформерным архитектурам, выпрямленным потокам, методам PEFT и переносу стиля. Для наглядного сопоставления рассмотренных методов приводится сравнительная таблица~\ref{tab:style_transfer_comparison}.

\begin{landscape}
\begin{table}[p]
\caption{Сравнение методов переноса стиля}
\label{tab:style_transfer_comparison}
\centering
\footnotesize
\begin{tabular}{p{2.8cm} p{2.2cm} p{4.5cm} p{2.5cm} p{2.0cm} p{2.5cm} p{2.0cm}}
\toprule
\textbf{Метод} & \textbf{Базовая модель} & \textbf{Подход} & \textbf{Обучение} & \textbf{Вход стиля} & \textbf{Сохранение содержания} & \textbf{Лицензия} \\
\midrule
Gatys et al.~\cite{gatys2016nst} (2016)
  & VGG
  & Итеративная оптимизация (матрица Грама)
  & Нет (on-image)
  & Изображение
  & Явный content loss
  & Research \\
\addlinespace
Textual Inversion~\cite{gal2023textualinversion} (2023)
  & SD 1.x
  & Оптимизация токен-эмбеддинга
  & Да (per-concept)
  & 3--5 изображений
  & Неявное через промпт
  & CreativeML \\
\addlinespace
DreamBooth~\cite{ruiz2023dreambooth} (2023)
  & SD 1.x/2.x
  & Полное дообучение + prior loss
  & Да (per-subject)
  & 3--5 изображений
  & Prior preservation
  & CreativeML \\
\addlinespace
IP-Adapter~\cite{ye2023ipadapter} (2023)
  & SDXL
  & Разделённый cross-attention адаптер
  & Да (общее)
  & Изображение
  & Сильное
  & Apache 2.0 \\
\addlinespace
StyleAligned~\cite{hertz2024stylealigned} (2024)
  & SDXL
  & Общие attention-статистики + AdaIN
  & Нет
  & Изображение
  & Умеренное
  & Research \\
\addlinespace
InstantStyle~\cite{wang2024instantstyle} (2024)
  & SDXL
  & Блочная инъекция признаков стиля
  & Нет
  & Изображение
  & Сильное
  & Apache 2.0 \\
\addlinespace
DEADiff~\cite{qi2024deadiff} (2024)
  & SDXL
  & Q-Former + нереконструктивное обучение
  & Да (общее)
  & Изображение
  & Сильное
  & Research \\
\addlinespace
\textbf{B-LoRA}~\cite{frenkel2024blora} \textbf{(2024)}
  & \textbf{SDXL}
  & \textbf{Поблочная LoRA, 2 блока, 1 фото}
  & \textbf{Да (${\approx}$5 мин/стиль)}
  & \textbf{1 изображение}
  & \textbf{Сильное}
  & \textbf{Research} \\
\addlinespace
IC-LoRA~\cite{huang2024iclora} (2024)
  & FLUX.1-dev
  & In-context LoRA
  & Да (малый датасет)
  & Набор изображений
  & Задачезависимое
  & Некоммерческая \\
\addlinespace
FluxSpace~\cite{dalva2024fluxspace} (2025)
  & FLUX.1-dev
  & Редактирование в attention-пространстве
  & Нет
  & Изображение
  & Умеренное
  & Research \\
\addlinespace
SplitFlux~\cite{yang2025splitflux} (2025)
  & FLUX.1-dev
  & Поблочная LoRA + RCA + VGRA
  & Да (${\approx}$5 мин/стиль)
  & 1 изображение
  & Сильное (RCA)
  & Research \\
\midrule
\textbf{B-LoRA-FLUX} (предлагаемый)
  & \textbf{FLUX.1-dev}
  & \textbf{Поблочная LoRA на MM-DiT}
  & \textbf{Да (${\approx}$5 мин/стиль)}
  & \textbf{1 изображение}
  & \textbf{Оценивается}
  & \textbf{---} \\
\bottomrule
\end{tabular}
\end{table}
\end{landscape}

Из изложенного складываются четыре магистральные направления развития рассматриваемой области. Во-первых, переход от пиксельных диффузионных моделей к латентным, а затем к архитектурам на основе трансформеров (DiT, MM-DiT, FLUX.1) перестроил организацию информационных потоков в генеративных моделях. Во-вторых, формализм выпрямленных потоков даёт ощутимые практические преимущества над классическими диффузионными процессами. В-третьих, методы PEFT, и LoRA в первую очередь, обеспечивают эффективную адаптацию больших моделей при минимальных вычислительных затратах. В-четвёртых, развитие переноса стиля идёт в сторону методов, опирающихся на структурную специализацию блоков модели.

Анализ выявил четыре конкретные лакуны, обосновывающие постановку задачи настоящего исследования: отсутствие прямого применения идей B-LoRA к FLUX.1; отсутствие систематического изучения специализации блоков FLUX.1; неисследованность роли блоков двойного потока в кодировании стилистической информации; отсутствие верификации применимости гипотезы блочной специализации при переходе от U-Net к DiT-архитектуре.

Перечисленные лакуны определяют цель и методологию работы. Раздел~\ref{sec:chapter3} предлагает методологические решения, адресующие эти пробелы: разработка и обоснование процедуры поиска специализированных блоков в FLUX.1, адаптация обучающей процедуры B-LoRA к трансформерной архитектуре и проектирование схемы инференса для независимого и совместного применения блочно-специфических адаптеров.
